% 65-37(65쪽) — 삼차함수 y=f(x) 의 그래프(최고차항 음수 · 극소는 y축 왼쪽 · 극대는 y축 바로 왼쪽 · f(0)>0).
% 원문에 있는 것만: 곡선 · 라벨 y=f(x), O, x, y. 눈금·좌표는 원문에 없으므로 넣지 않는다(계수의 부호를 묻는 문항이라 수치를 더하면 답이 드러난다).
% 스캔 화질이 낮아 TikZ 로 다시 그림.
\begin{tikzpicture}[x=0.80cm, y=0.52cm, line width=0.5pt]
  \draw[->, >=stealth, line width=0.7pt] (-3.00,0) -- (1.30,0) node[below=1pt, font=\small] {$x$};
  \draw[->, >=stealth, line width=0.7pt] (0,-1.75) -- (0,2.70) node[left=1pt, font=\small] {$y$};
  % 2026-10-02 검수: 원본 x절편 비 -3.61 : -1.39 : 1 에 맞춰 근을 -1.63, -0.63, 0.45 로 다시 잡았다.
  \draw[line width=0.6pt, domain=-2.05:0.95, samples=120, smooth]
    plot (\x, {-0.8*(((\x)+1.63)*((\x)+0.63)*((\x)-0.45))});
  \node[below right=-1pt and 0pt, font=\small] at (0,0) {$\pt{O}$};
  \node[anchor=east, font=\small] at (-1.55,2.25) {$y=f(x)$};
\end{tikzpicture}
